% =============================================================
%  Chapter 2 — Related Work (full draft)
% =============================================================
%  Matches the structure in 01-paper-analysis.md:
%    §2.1  LLM code generation and backend evaluation benchmarks
%    §2.2  Reliable and verifiable LLM code generation
%    §2.3  Deployment configuration and cluster tuning
%    §2.4  LLM-driven performance optimization
%    §2.5  Positioning of this work
% =============================================================

\chapter{Related Work}
\label{ch:related-work}


% ----- §2.1  Benchmarks -----
\section{LLM Code Generation and Backend Evaluation Benchmarks}
\label{sec:rw-benchmarks}

Evaluation of large language models for software engineering has progressed
along a clear trajectory of increasing scope and realism.
HumanEval~\cite{chen2021humaneval} established pass@$k$ as the standard metric
for single-function Python synthesis from natural-language specifications.
MBPP and related suites extended this to short programs with broader test
coverage.  SWE-bench~\cite{jimenez2024swebench} moved evaluation to the
repository level, requiring models to resolve real GitHub issues by producing
patches that pass project test suites.  These benchmarks measure whether
generated software is \emph{correct} at increasing scales of context, but they
do not assess complete backend services, security properties, or runtime
performance under deployment.

\textsc{BaxBench}~\cite{vero2025baxbench} addresses the application layer:
392~tasks across 28~scenarios, 14~frameworks, and 6~languages, evaluating
whether LLMs can generate \emph{complete} backend applications that pass
functional tests and resist security exploits.  Results show that even strong
models achieve only moderate functional correctness and that a large fraction of
correct code remains insecure.  Complementary work by CWEval~\cite{peng2025cweval}
jointly evaluates functionality and security with outcome-driven test oracles,
demonstrating that prior security benchmarks can substantially overestimate
safety.  SecRepoBench~\cite{shen2026secrepobench} extends secure-code evaluation
to repository-level \emph{code agents}, combining developer unit tests with
OSS-Fuzz security tests across 318~completion tasks in real C/C++ projects.

CodeFlowBench~\cite{wang2025codeflowbench} introduces a complementary dimension:
\emph{multi-turn iterative} code generation (``codeflow''), where new
functionality is built by reusing functions implemented in earlier turns.
Experiments show significant performance degradation relative to single-turn
generation as dependency complexity grows, highlighting that iterative software
construction is a distinct and difficult evaluation setting.

Together, these benchmarks characterise whether LLM-generated software is
correct, secure, or constructible over multiple turns.  However, they
typically execute programs in isolated containers without orchestrated
deployment, shared cluster resources, or load testing.  None measures
\emph{sustainable throughput} when an application is deployed as a replicated
service on Kubernetes---the gap our work addresses.


% ----- §2.2  Reliability -----
\section{Reliable and Verifiable LLM Code Generation}
\label{sec:rw-reliability}

Because unconstrained LLM output is often syntactically invalid or semantically
incorrect, a line of work seeks to improve reliability before any deployment
attempt.  Synchromesh~\cite{poesia2022synchromesh} combines Target Similarity
Tuning for few-shot example selection with Constrained Semantic Decoding, which
restricts token sampling to programs satisfying grammar, typing, and domain
constraints---eliminating whole classes of parse and runtime errors without
fine-tuning the underlying language model.

Clover~\cite{sun2023clover} proposes \emph{closed-loop verifiable} code
generation: the model produces code together with formal annotations and
natural-language docstrings, and a consistency checker validates alignment among
all three using formal verification tools coupled with LLMs.  On textbook-level
Dafny programs, the approach achieves high acceptance of correct instances while
maintaining zero tolerance for adversarial incorrect ones.

VERGE~\cite{singh2026verge} applies a related closed-loop pattern to logical
reasoning: LLM outputs are decomposed into claims, autoformalized, and verified
with SMT solvers; Minimal Correction Subsets provide actionable refinement
feedback across iterations.

Our framework shares the principle that LLM artifacts must pass hard gates
before downstream stages proceed, but we enforce reliability through BaxBench
functional tests and schema-validated deployment specifications rather than
formal verification or constrained decoding of source code.  Reliability in our
setting is necessary yet insufficient: a deployment that passes tests may still
perform poorly under load, motivating the performance-oriented loop described in
Sections~\ref{sec:rw-cluster} and~\ref{sec:rw-llm-perf}.


% ----- §2.3  Cluster tuning -----
\section{Deployment Configuration and Cluster Tuning}
\label{sec:rw-cluster}

Production backends run on container orchestration platforms rather than in
isolation.  Burns~et~al.~\cite{burns2016borg} document how Google's Borg,
Omega, and Kubernetes evolved to schedule workloads, enforce resource isolation,
and maintain declarative desired state across large clusters.  In Kubernetes,
operators express deployment intent through replica counts, CPU and memory
requests and limits, affinity rules, and service definitions; the control plane
continuously reconciles cluster state with that intent.

Day-to-day performance tuning often relies on \emph{reactive} autoscaling.
The Horizontal Pod Autoscaler (HPA)~\cite{kubernetes-hpa} adjusts replica
count from live metrics such as CPU utilisation or custom request rates; the
Vertical Pod Autoscaler adjusts per-pod resource requests.  These mechanisms
optimise scale within a \emph{fixed} deployment template.  They do not reason
about database topology, connection pooling, pod placement across worker nodes,
or joint trade-offs between application code and infrastructure layout.  Nor do
they perform multi-iteration search guided by structured post-load-test
diagnostics.

Learned cluster management predates LLM agents.  Decima~\cite{mao2019decima}
uses deep reinforcement learning to discover scheduling policies for data-
processing clusters that outperform hand-tuned heuristics.  Like classical
knob-tuning systems, Decima optimises \emph{policies or parameters} within a
fixed system architecture rather than revising application code or declarative
deployment manifests for arbitrary HTTP services.

More recently, Glia~\cite{hamadanian2025glia} applies a multi-agent LLM
architecture---combining reasoning, experimentation, and analysis---to design
systems mechanisms for a distributed GPU cluster running LLM inference.  The
system autonomously proposes algorithms for request routing, job scheduling,
and autoscaling that match human-expert quality in substantially less time, with
interpretable designs grounded in empirical feedback.  Glia is the closest prior
work to our setting in combining \emph{LLM reasoning} with \emph{measured
performance feedback} for infrastructure behaviour.  Our work differs in three
respects: we target \emph{general BaxBench backend applications} rather than
inference-serving stacks; we optimise a \emph{constrained deployment
specification} (\texttt{spec.yaml}) rather than generating new routing or
scheduling algorithms; and we score iterations by \emph{adaptive load-test
goodput} on a multi-node Kubernetes cluster with explicit correctness gates
before each benchmark.

No prior system in this line evaluates whether an LLM agent can \emph{iteratively}
refine replica topology, resource allocation, placement, and optional
application code for LLM-generated web backends using sustainable-throughput
feedback---the problem our framework formalises.


% ----- §2.4  Performance optimization -----
\section{LLM-Driven Performance Optimization}
\label{sec:rw-llm-perf}

A growing body of work asks whether LLMs can improve \emph{performance}, not
merely correctness.  SWE-Perf~\cite{he2025sweperf} is the first benchmark for
repository-level \emph{code} performance optimisation: 140~instances derived
from performance-improving pull requests, each with performance tests and
expert reference patches.  Evaluations of Agentless~\cite{xia2024agentless},
OpenHands, and related approaches reveal a substantial gap between current LLMs
and expert optimisers.  The benchmark is largely \emph{one-shot}: the model
receives a codebase and must emit an optimised patch without iterative runtime
feedback from executing the system under load.

Agentic systems have begun to close the loop with measurement.
VibeServe~\cite{kamahori2025vibeserve} synthesises bespoke LLM serving runtimes
through a two-level multi-agent loop.  An outer planner maintains an issue
tracker and git-recorded optimisation history; an inner loop separates
implementation, accuracy judging, and performance evaluation.  On standard
Llama serving, generated systems match vLLM; on non-standard model--hardware--
workload combinations, they achieve large speedups by specialising kernels,
memory management, and scheduling logic.

Liu~et~al.~\cite{liu2025jitsystems} argue for \emph{Just-in-Time Systems}:
LLM agents can synthesise entire key-value stores from specification cards
describing environment, workload, and required properties.  Their pipeline
(\textsc{Jitskit}) iterates through planning, coding, evaluation with leading
indicators, and critic-guided reflection, with an adversarial auditor that
detects reward hacking.  Synthesised stores outperform FASTER, RocksDB, and
Redis on all 18~configurations tested (up to $4.6\times$ on the most favourable
spec).

Our framework shares the iterative structure of VibeServe and Jitskit---each
round produces a candidate, passes correctness gates, is benchmarked, and
feeds structured diagnostics to the next iteration---but operates at a different
level of the stack.  SWE-Perf optimises \emph{source code}; VibeServe and
Jitskit synthesise \emph{system implementations}; we optimise \emph{Kubernetes
deployment configuration} (and optionally application code) for BaxBench
backends, using \emph{goodput}---sustained successful requests per second---as
the primary objective under an adaptive load profile.


% ----- §2.5  Positioning -----
\section{Positioning of This Work}
\label{sec:rw-positioning}

Table~\ref{tab:rw-comparison} summarises how representative systems relate to
our framework across evaluation target, feedback structure, and deployment
environment.  Benchmarks through CodeFlowBench establish correctness, security,
and multi-turn code quality but not cluster performance.  Reliability methods
(Synchromesh, Clover) and cluster tooling (autoscaling, Decima, Glia) address
trust and infrastructure tuning without benchmarking full LLM-generated backends
under adaptive load.  SWE-Perf, VibeServe, and Jitskit demonstrate LLM-driven
performance optimisation at the code or implementation level.

Our thesis occupies the intersection BaxBench leaves open: \emph{iterative
deployment optimisation} of complete backend applications on a real Kubernetes
cluster, with the agent controlling both application code and a validated
deployment specification, and with each iteration scored by measured goodput
rather than static tests alone.

\begin{table}[t]
  \centering
  \caption{Comparison of related approaches (selected dimensions).}
  \label{tab:rw-comparison}
  \small
  \begin{tabular}{@{}l c c c c c@{}}
    \toprule
    & \textbf{BaxBench} & \textbf{Glia} & \textbf{SWE-Perf} &
      \textbf{VibeServe} & \textbf{This work} \\
    \midrule
    Primary focus
      & Func.+sec.
      & Mechanism design
      & Code perf.
      & Serving runtime
      & Deploy perf. \\
    Iterative loop
      & No
      & Yes
      & No
      & Yes
      & Yes \\
    Runtime load test
      & No
      & Yes
      & Perf.\ tests
      & Yes
      & Yes \\
    Environment
      & Docker
      & GPU cluster
      & Local/CI
      & Single GPU
      & Multi-node K8s \\
    Agent controls
      & Code
      & Algorithms
      & Code patches
      & System code
      & Code + spec \\
    \bottomrule
  \end{tabular}
\end{table}
